% Demonstrator: linguexx examples in a tagged (PDF/UA-oriented) document.
% Compile with lualatex (recommended for tagging), pdflatex, or xelatex.
% The \DocumentMetadata line is what turns tagging on; remove it and the
% same source typesets identically, untagged.
% Tagging is enabled by the \DocumentMetadata line below.
% On TeX Live 2025-11 or newer, \DocumentMetadata{lang=en,tagging=on} is
% the recommended form.  testphase={phase-III} works on TeX Live 2023
% through 2026 and selects the current comprehensive tagging.
\DocumentMetadata{lang=en,pdfversion=2.0,testphase={phase-III}}
\documentclass{article}
\usepackage[lazy,gb4e]{linguexx}
\GlossTierLang{1}{de} % object line is German
\begin{document}

\section{Tagged linguistic examples}

Ordinary running text precedes the first example.

\ex.\label{intro} A numbered example, tagged as a labelled list item.
\a. A sub-example at the letter level.
\b. *A judged sub-example.
\a. A sub-sub-example at the roman level.
\z.\z.

A gloss, whose tiers are word-aligned:

\exg. Der Hund bellte laut.\\
      the dog bark.\lpzg{pst} loudly.\\
\glt `The dog barked loudly.'

Examples survive inside footnotes,\footnote{Footnote text with its own
example: \ex. A footnoted example. }
straddle page boundaries, and interleave with cross-references such as
\ref{intro}.

\end{document}
